\renewcommand{\Anno}{2009}

\renewcommand{\NumeroProblema}{2}

\renewcommand{\Testo}{
Per cercare di arginare il download di musica illegale le major discografiche hanno messo a punto un pericoloso virus informatico: una volta installatosi su un computer, quando l'utente tenta di scaricare un brano musicale, il virus con probabilità $p$ (con $0<p<1$) permette lo scaricamento senza alcuna conseguenza, altrimenti lo scaricamento non avviene, compare una scritta minacciosa ed il virus memorizza un ``tag''.  Ovviamente l'utente può ritentare lo scaricamento, ma al terzo tag memorizzato (anche su tentativi effettuati per brani diversi) il virus formatta l'hard disk.

Tullio vorrebbe scaricare dei brani musicali, ma il suo computer è stato infettato dal virus.
\begin{enumerate}
    \item Determinare, in funzione di $p$, la probabilità che Tullio riesca a scaricare 2 brani in \emph{esattamente} 4 tentativi.

    \item  Determinare, in funzione di $p$, la probabilità che Tullio riesca a scaricare 2 brani.

    \item Determinare, in funzione di $p$ ed $n$, la probabilità che Tullio riesca a scaricare $n$ brani.
\end{enumerate}
}

\renewcommand{\Risposta}{\begin{enumerate}
		\item[(a)]   $3p^{2}(1-p)^{2}$
	
		\item[(b)]  $p^{2}(3p^{2}-8p+6)$
		
		\item[(c)] $\displaystyle{p^{n}+np^{n}(1-p)+\frac{n(n+1)}{2}p^{n}(1-p)^{2}}$
	\end{enumerate}}

\renewcommand{\Aiutino}{\begin{enumerate}
	\renewcommand{\labelenumi}{(\alph{enumi})} 
	\item[(a)]  Quanti tentativi non sono andati a buon fine? In quale posizione possono trovarsi questi tentativi?

	\item[(b)] Si tratta di calcolare e poi sommare le probabilità che Tullio ce la faccia in 2, 3 o 4 tentativi.
	
	\item[(c)] Si tratta di calcolare e poi sommare le probabilità che Tullio ce la faccia in $n$, $(n+1)$ o $(n+2)$ tentativi. Prestare particolare attenzione alle possibili posizioni del tentativo (o dei tentativi) andati male.
\end{enumerate}}

\renewcommand{\SoluzioneUfficiale}{\TitoloSoluzione\DomandaA
La probabilità richiesta è $3p^{2}(1-p)^{2}$.

Se Tullio riesce a scaricare 2 brani in esattamente 4 tentativi, vuol dire che esattamente 2 dei primi 3 tentativi non sono andati a buon fine, con conseguente memorizzazione del tag. Abbiamo quindi 3 possibili sequenze che realizzano la situazione proposta, e precisamente
$$\mbox{ok--no--no--ok},
\hspace{3em}
\mbox{no--ok--no--ok},
\hspace{3em}
\mbox{no--no--ok--ok}.$$

Si noti in particolare che il quarto tentativo è andato necessariamente bene. Ciascuna delle 3 sequenze ha probabilità $p^{2}(1-p)^{2}$, in quanto contiene 2 successi e 2 insuccessi. Trattandosi di eventi disgiunti, non resta che sommare le probabilità, ottenendo appunto $3p^{2}(1-p)^{2}$.

\DomandaB 
La probabilità richiesta è
$$p^{2}+2p^{2}(1-p)+3p^{2}(1-p)^{2}=p^{2}(3p^{2}-8p+6).$$

Il primo addendo, $p^{2}$, rappresenta la probabilità che Tullio riesca a scaricare i 2 brani in soli 2 tentativi (2 successi consecutivi).

Il secondo addendo, $2p^{2}(1-p)$, rappresenta la probabilità che Tullio riesca a scaricare i 2 brani in \emph{esattamente} 3 tentativi. In tal caso infatti sono possibili 2 situazioni disgiunte (ok--no--ok e no--ok--ok), le quali hanno entrambe probabilità $p^{2}(1-p)$ perché contengono 2 successi ed un mancato scaricamento.

Il terzo addendo, $3p^{2}(1-p)^{2}$, rappresenta la probabilità che Tullio riesca a scaricare i 2 brani in \emph{esattamente} 4 tentativi.  Tale probabilità è stata calcolata precedentemente.

È chiaro infine che Tullio non può impiegare più di 4 tentativi per scaricare 2 brani, in quanto questo comporterebbe la memorizzazione di 3 o più tag, con conseguente formattazione dell'hard disk.


\DomandaC 
La probabilità richiesta è
$$p^{n}+np^{n}(1-p)+\frac{n(n+1)}{2}p^{n}(1-p)^{2}.$$

Per dimostrarlo osserviamo preliminarmente che lo scaricamento di $n$ brani può avvenire soltanto in $n$, $n+1$, o $n+2$ tentativi.

La probabilità che Tullio ce la faccia in $n$ tentativi è $p^{n}$, perché devono essere tutti successi.

La probabilità che Tullio ce la faccia in \emph{esattamente}
$(n+1)$ tentativi è $np^{n}(1-p)$.  Si devono infatti realizzare $n$ successi ed un insuccesso, il quale può situarsi in uno qualunque dei primi $n$ tentativi (ma non nell'ultimo!).

Calcoliamo infine la probabilità che Tullio ce la faccia in \emph{esattamente} $(n+2)$ tentativi. In tal caso servono $n$ successi e 2 insuccessi, i quali possono situarsi in 2 qualsiasi dei primi $(n+1)$ tentativi (ma non nell'ultimo!). I modi di scegliere 2 tentativi infruttuosi su $(n+1)$ sono $\binom{n+1}{2}$. Ne segue che quest'ultima probabilità è data da
$$\binom{n+1}{2}p^{n}(1-p)^{2}=\frac{n(n+1)}{2}p^{n}(1-p)^{2}.$$
}

\renewcommand{\SoluzioniAlternative}{}

\renewcommand{\Osservazioni}{\TitoloOsservazione
Può essere interessante considerare il caso più generale in cui la formattazione dell'hard disk avviene al $k$-esimo tag. Ragionando come nella soluzione proposta si ottiene che la probabilità di scaricare $n$ brani è in tal caso data dalla formula
$$p^{n}\sum_{i=0}^{k-1}\binom{n+i-1}{i}(1-p)^{i}.$$

Per $k=3$ si riottiene ovviamente quanto dimostrato precedentemente.
}

\renewcommand{\Commenti}{
}